
\chapter{攻读硕士学位期间取得的研究成果}

{
\zihao{5}
% \setstretch{1.65}
\vspace{2pt}
一、已发表（包括已接受待发表）的论文，以及已投稿、或已成文打算投稿、或拟成文投稿的论文情况\underline{\textbf{（只填写与学位论文内容相关的部分）：}}
\vspace{1em}
% ==================== 表格区域 ====================
% 定义列格式：m代表垂直居中，后面的长度是列宽，<{\centering}代表水平居中
\begin{table}[h]
    \centering
    \zihao{-5}
    \renewcommand{\arraystretch}{3}
    \renewcommand{\tabularxcolumn}[1]{m{#1}}
    \begin{tabularx}{\textwidth}{|m{0.6cm}<{\centering}|X<{\centering}|m{2.0cm}<{\centering}|m{2.0cm}<{\centering}|m{2.4cm}<{\centering}|m{2.0cm}<{\centering}|}
        \hline
        序号 & 
        发表或投稿刊物/会议名称 & 
        作者（仅注明第几作者） & 
        发表年份 & 
        与学位论文哪一部分（章、节）相关 & 
        被索引收录情况 \\ 
        \hline
        
        % 第一行数据
        1 & 
        IEEE/CVF Computer Vision and Pattern Recognition Conference \textbf{(CVPR)} & 
        第一作者 &
        2026（已接受） & 
        第三章 & 
        无 \\ 
        \hline
        
    \end{tabularx}
\end{table}

\vspace{0.5em}

% % ==================== 注释区域 ====================
% \noindent 注：在“发表的卷期、年月、页码”栏：

% \noindent 1. 如果论文已发表，请填写发表的卷期、年月、页码；

% \noindent 2. 如果论文已被接受，填写将要发表的卷期、年月；

% \noindent 3. 以上都不是，请据实填写“已投稿”，“拟投稿”。

% \noindent 不够请另加页。

% \vspace{1.5em}

% % ==================== 第二部分 ====================
% \noindent 二、与学位内容相关的其它成果（包括专利、著作、获奖项目等）

% \vspace{0.5em}

% \noindent 专利成果：

% \begin{itemize}
%     \item 杜卿、\textbf{刘东}、杨逸凡、谭明奎. 基于高频局部特征的三维穿衣人体重建方法、装置及介质. 
%     CN202311764613.8.
% \end{itemize}
% }